
\documentclass[a4paper,12pt]{article}

\usepackage{JYM-harj}
\usepackage{tikz}
\usepackage{amsmath}
\usepackage{tikz}

\setlength{\parindent}{0in}

\setlength{\voffset}{-1.0cm}
\addtolength{\textheight}{2.0cm}

\setlength{\hoffset}{-1.0cm}
\addtolength{\textwidth}{2.0cm}

\usepackage{graphicx}
\usepackage{verbatim}

\usepackage{hyperref}

%\usepackage[finnish]{babel}
%\usepackage[utf8]{inputenc}
%\usepackage[T1]{fontenc}
%\usepackage{amsmath,amsfonts,amssymb,amsthm,amscd}

\newcommand{\tversio}{1}

\newcommand{\harjoitusnro}{6}

\ifthenelse{\tversio = 1}{\pagestyle{empty}}{}

\begin{document} 

\otsikko

\begin{enumerate}

\item Ratkaise yht\"al\"o $z^6=64$ kompleksilukujen joukossa. Piirr\"a ratkaisut kompleksitasoon.

\textbf{Vastaus:} Ratkaistaan $z^6=64$:

\begin{align*}
    64&=64e^{0i}\\
    z^6&=(re^{\phi i})^6 =r^6e^{6\phi i}\\    
\end{align*}
\[
\begin{cases}
    r^6=64 \\
    6\phi = k \cdot 2\pi
\end{cases}
\quad \Longleftrightarrow \quad
\begin{cases}
    r=\sqrt[6]{64} \\
    \phi=k\cdot \frac{2}{6}\cdot \pi 
\end{cases}
\quad \Longleftrightarrow \quad
\begin{cases}
    r=2 \\
    \phi=k\cdot \frac{1}{3}\cdot \pi 
\end{cases}
\]

Ratkaisut ovat: $2e^{0i}$, $2e^{\frac{1}{3}\pi i}$, $2e^{\frac{2}{3}\pi i}$, $2e^{\pi i}$, $2e^{\frac{4}{3}\pi i}$, $2e^{\frac{5}{3}\pi i}$. Ne on piirretty alla olevaan kuvaan.\\

\begin{center}
\begin{tikzpicture}[scale=1.5]
    \draw[->] (-3,0) -- (3,0) node[right] {Re};
    \draw[->] (0,-3) -- (0,3) node[above] {Im};

    \foreach \x in {-2, -1, 1, 2} {
        \draw (\x, 0.1) -- (\x, -0.1);
    }

    \foreach \y in {-2, -1, 1, 2} {
        \draw (0.1, \y) -- (-0.1, \y);
    }

    \draw[dashed, gray] (0,0) circle (2);

    \foreach \angle in {0, 60, ..., 300} {
        \filldraw[blue] (\angle:2) circle (2pt);
        %\draw (0,0) -- (\angle:2);
    }
    
    \node[below right] at (2,0) {$2$};
    \node[below left] at (-2,0) {$-2$};

    \node[above left] at (0,2) {$2i$};
    \node[below left] at (0,-2) {$-2i$};


    \node[align=center, below] at (0,-3.5) {Yhtälön $z^6=64$ ratkaisut};
    
\end{tikzpicture}
\end{center}

\item Ratkaise yht\"al\"o $(z-4)^4=81$ kompleksilukujen joukossa. Piirr\"a ratkaisut kompleksitasoon.

\textbf{Vastaus:} Ratkaistaan $(z-4)^4=81$. Otetaan yhtälön vasemmalle puolelle apumuuttuja $w$, jonka avulla saadaan $w=(z-4)$, jolloin myös $w^4=(z-4)^4$ pätee. Jatketaan eksponenttiesityksellä:
\begin{align*}
    81&=81e^{0i}\\
    w^4&=r^4e^{4\phi i}
\end{align*}
\[
\begin{cases}
    r^4=81\\
    4\phi = k\cdot 2\pi
\end{cases}
\quad \Longleftrightarrow \quad
\begin{cases}
    r=\sqrt[4]{81}\\
    \phi=k\cdot \frac{2}{4}\cdot \pi
\end{cases}
\quad \Longleftrightarrow \quad
\begin{cases}
    r=3\\
    \phi=k\cdot \frac{1}{2}\cdot \pi
\end{cases}
\]

Kun $k=0$, saadaan $w_0=3e^{0i}=3$, kun $k=1$, saadaan $w_1=3e^{\frac{1}{2}\pi i}=3i$, kun $k=2$, saadaan $w_2=3e^{\pi i}=-3$, ja kun $k=3$, saadaan $w_3=3e^{\frac{3}{2}\pi i}=-3i$. 

Koska $w=z-4$, voidaan muuttujan $z$ arvot ratkaista: $z=w+4$. Näin saadaan $z_0=3+4=7$, $z_1=3i+4=4+3i$, $z_2=-3+4=1$ ja $z_3=-3i+4=4-3i$. Nämä pisteet on piirretty alla olevaan kuvaan:

\begin{center}
\begin{tikzpicture}[scale=1.1]
    \draw[->] (-2,0) -- (10,0) node[right] {Re};
    \draw[->] (0,-4) -- (0,4) node[above] {Im};

    \foreach \x in {-1, 1, 2, 3, 4, 5, 6, 7, 8, 9} {
        \draw (\x, 0.1) -- (\x, -0.1);
    }

    \foreach \y in {-3, -2, -1, 1, 2, 3} {
        \draw (0.1, \y) -- (-0.1, \y);
    }

    \draw[dashed, gray!50] (4,0) circle (3);
    \draw[dashed, gray] (7,0) -- (4,3) -- (1,0) -- (4,-3) -- cycle;

    \filldraw[blue] (7,0) circle (2pt);
    \filldraw[blue] (4,3) circle (2pt);
    \filldraw[blue] (1,0) circle (2pt);
    \filldraw[blue] (4,-3) circle (2pt);

    \node[below] at (1,-0.1) {$1$};
    \node[below] at (4,-0.1) {$4$};
    \node[below] at (7,-0.1) {$7$};
    \node[below] at (10,-0.1) {$10$};

    \node[left] at (-0.1, 3) {$3i$};
    \node[left] at (-0.1, -3) {$-3i$};
    \node[below left] at (0,0) {$0$};

    \node[align=center, below] at (5,-4.5) {Yhtälön $(z-4)^4=81$ ratkaisut};
\end{tikzpicture}
\end{center}
\vspace{1em}

\item Piirr\"a kompleksitasoon $|z|=3$ ja $|z-3|=4$. Selvit\"a pisteet, joissa kuviot leikkaavat k\"aytt\"am\"att\"a laskinohjelmistoja.

\textbf{Vastaus:}

\begin{center}
\begin{tikzpicture}[scale=0.8]
    \draw[->] (-4,0) -- (10,0) node[right] {Re};
    \draw[->] (0,-5) -- (0,5) node[above] {Im};

    \foreach \x in {-3, -2, -1, 1, 2, 3, 4, 5, 6, 7, 8, 9} {
        \draw (\x, 0.1) -- (\x, -0.1);
    }

    \foreach \y in {-4, -3, -2, -1, 1, 2, 3, 4} {
        \draw (0.1, \y) -- (-0.1, \y);
    }

    \draw[dashed, blue] (0,0) circle (3);
    \draw[dashed, red] (3,0) circle (4);

    \filldraw[blue] (0,0) circle (2pt);
    \filldraw[red] (3,0) circle (2pt);

    \node[below left] at (-3,0) {$-3$};
    \node[below right] at (3,-0.1) {$3$};
    \node[below] at (7,-0.1) {$7$};
    \node[below] at (10,-0.1) {$10$};

    \node[left] at (-0.1, 4) {$4i$};
    \node[left] at (-0.1, -4) {$-4i$};
    \node[below left] at (0,0) {$0$};

    \node[align=center, below] at (3,-5.5) {Kuvassa sinisellä värillä  $|z|=3$ ja punaisella värillä $|z-3|=4$.};
\end{tikzpicture}
\end{center}

Selvitetään sitten tarkat pisteet, joissa kuviot leikkaavat. Meillä on siis kaksi muuttujaa, joissa molemmissa on $z$. Muuttuja $z$ voidaan ilmaista myös muodossa $x+yi$. Sijoitetaan se yhtälöihin muuttujan $z$ paikalle. Näin saadaan: $|x+yi|=3$ ja $|(x-3)+yi|=4$. Näistä saadaan Pythagoraan lauseella yhtälöpari:
\begin{itemize}
    \item $|x+yi|=3 \quad \Longrightarrow \quad x^2+y^2=9$
    \item $|(x-3)+yi|=4\quad \Longrightarrow \quad (x-3)^2+y^2=16$
\end{itemize}

Ratkotaan $x$ jälkimmäisestä yhtälöstä:
\begin{align*}
    (x-3)^2+y^2&=16\\
    x^2+2\cdot x\cdot (-3)+(-3)^2+y^2&=16\\
    x^2-6x+9+y^2&=16\\
    -6x+9+9&=16 \quad \text{(sijoitetaan toinen yhtälö)}\\
    18-6x&=16\\
    -6x&=-2\\
    x&=\frac{1}{3}
\end{align*}

Sijoitetaan muuttujan $x$ arvo yhtälöön $x^2+y^2=9$:
\begin{align*}
    \left(\frac{1}{3}\right)^2+y^2&=9\\
    y^2&=9-\frac{1}{9} \\
    y&= \pm\sqrt{\frac{80}{9}}\\
    y&=\pm \frac{4\cdot\sqrt{5}}{3}
\end{align*}

Nyt kun meillä on $x$ ja $y$, voidaan johtaa kuvioiden leikkaupisteet: 
\begin{itemize}
    \item $\left( \frac{1}{3},\frac{4\cdot\sqrt{5}}{3}i\right) \quad \approx \quad (0.33,2.98i)$
    \item $\left(\frac{1}{3},-\frac{4\cdot\sqrt{5}}{3}i\right) \quad \approx \quad (0.33,-2.98i)$.
\end{itemize}
\vspace{1em}

\item Todista eksponenttimuodon avulla trigonometrian kaava
\[
\cos(2\varphi)=\cos^2\varphi-\sin^2\varphi.
\]

\textbf{Vastaus:} Tarkastellaan lauseketta $e^{i2\varphi}$ Eulerin kaavan $e^{i\varphi}=\cos\varphi+i\sin\varphi$ avulla. Saadaan lauseke 
\begin{equation*}
    e^{i2\varphi}=\cos(2\varphi)+i\sin(2\varphi).
\end{equation*}

Tarkastellaan lukua $e^{i2\varphi}$ toisella tavalla:
\begin{equation*}
    e^{i2\varphi}=\left(e^{i\varphi}\right)^2=(\cos\varphi+i\sin\varphi)^2.
\end{equation*}
Avataan jälkimmäinen lauseke:
\begin{align*}
    (\cos\varphi+i\sin\varphi)^2 &= (\cos\varphi)^2 + 2\cos\varphi(i\sin\varphi) + (i\sin\varphi)^2 \\
    &= \cos^2\varphi + i(2\sin\varphi\cos\varphi) + i^2\sin^2\varphi \\
    &= (\cos^2\varphi - \sin^2\varphi) + i(2\sin\varphi\cos\varphi).
\end{align*}
Koska meillä on nyt kaksi lauseketta, jotka molemmat kuvaavat samaa kompleksilukua $e^{i2\varphi}$, niiden täytyy olla yhtä suuret keskenään:
\begin{equation*}
    \cos(2\varphi)+i\sin(2\varphi)=(\cos^2\varphi- \sin^2\varphi)+i(2\sin\varphi\cos\varphi).
\end{equation*}
Nyt kun vertaamme näiden yhtälöiden reaaliosia, saadaan:
\begin{equation*}
    \cos(2\varphi)=\cos^2\varphi-\sin^2\varphi.
\end{equation*}
Kaava on näin todistettu oikeaksi.

\item Todista eksponenttimuodon avulla trigonometrian kaava
\[
2\sin \theta\cos \varphi=\sin(\theta+\varphi)+\sin(\theta-\varphi)
\]

\textbf{Vastaus:} Tarkastellaan kaavan oikeaa puolta $e^{i(\theta+\varphi)} + e^{i(\theta-\varphi)}$ kahdella eri tavalla. 

1: Avataan termit Eulerin kaavalla $e^{i\alpha} = \cos\alpha + i\sin\alpha$:
\begin{align*}
    e^{i(\theta+\varphi)} + e^{i(\theta-\varphi)} &= \big(\cos(\theta+\varphi) + i\sin(\theta+\varphi)\big) + \big(\cos(\theta-\varphi) + i\sin(\theta-\varphi)\big) \\
    &= \big(\cos(\theta+\varphi) + \cos(\theta-\varphi)\big) + i\big(\sin(\theta+\varphi) + \sin(\theta-\varphi)\big).
\end{align*}

2: Käytetään eksponenttien laskusääntöjä:
\begin{equation*}
    e^{i(\theta+\varphi)} + e^{i(\theta-\varphi)} = e^{i\theta}e^{i\varphi} + e^{i\theta}e^{-i\varphi} = e^{i\theta}\left(e^{i\varphi} + e^{-i\varphi}\right).
\end{equation*}
Sovelletaan sulkulausekkeeseen Eulerin kaavaa: 
\begin{align*}
    e^{i\varphi} + e^{-i\varphi} &= (\cos\varphi + i\sin\varphi) + (\cos(-\varphi) + i\sin(-\varphi)) \\
    &= \cos\varphi + i\sin\varphi + \cos\varphi - i\sin\varphi \\
    &= 2\cos\varphi.
\end{align*}
Sijoitetaan tämä takaisin lausekkeeseen $e^{i\theta}\left(e^{i\varphi} + e^{-i\varphi}\right)$:
\begin{align*}
    e^{i\theta}\left(e^{i\varphi} + e^{-i\varphi}\right) &= (\cos\theta + i\sin\theta)(2\cos\varphi) \\
    &= 2\cos\theta\cos\varphi + i(2\sin\theta\cos\varphi).
\end{align*}

Koska tavat 1 ja 2 kuvaavat samaa summalauseketta, niiden on oltava yhtä suuret keskenään:
\begin{equation*}
    \big(\cos(\theta+\varphi) + \cos(\theta-\varphi)\big) + i\big(\sin(\theta+\varphi) + \sin(\theta-\varphi)\big) = 2\cos\theta\cos\varphi + i(2\sin\theta\cos\varphi).
\end{equation*}

Kun vertaamme yhtälöiden imaginääriosia, saadaan:
\begin{equation*}
    \sin(\theta+\varphi) + \sin(\theta-\varphi) = 2\sin\theta\cos\varphi.
\end{equation*}
Kaava on näin todistettu oikeaksi.

\item Vastaa seuraaviin kysymyksiin: \begin{enumerate}
    \item Mitk\"a ovat polynomin $z^2+9$ nollakohdat kompleksilukujen joukossa?
    
    \textbf{Vastaus:} 
    \begin{align*}
        z^2+9&=0\\
        z^2&=-9\\
        z&=\pm i\sqrt{|-9|}\\
        z&=\pm3i
    \end{align*}
    Polynomin nollakohdat ovat $z=3i$ ja $z=-3i$. 
    
    \item Mitk\"a ovat polynomin $z^2+2z+2$ nollakohdat kompleksilukujen joukossa?

    \textbf{Vastaus:}
    \begin{align*}
        z^2+2z+2&=0\\
        z&=\frac{-2\pm \sqrt{2^2-4\cdot 1\cdot 2}}{2}\\
        &=\frac{-2\pm i \sqrt{|-4|}}{2}\\
        &=\frac{-2\pm i 2}{2}\\
        &=-1\pm i
    \end{align*}
    
    \item  Mitk\"a ovat polynomin $z^4-1$ nollakohdat kompleksilukujen joukossa?
    
    \textbf{Vastaus:}
    \begin{align*}
        z^4-1&=0\\
        (z^2-1)(z^2+1)&=0
    \end{align*}

    Tulon nollasäännön mukaisesti joko ensimmäinen tai toinen tulon tekijä on nolla. Eli joko
    \begin{align*}
        z^2-1&=0\\
        z^2&=1\\
        z=1 \quad &\text{TAI} \quad z=-1
    \end{align*}
    tai
    \begin{align*}
        z^2+1&=0\\
        z^2&=-1\\
        z=i \quad &\text{TAI} \quad z=-i
    \end{align*}

    Polynomilla on siis neljä nollakohtaa: $z=1$, $z=-1$, $z=i$ ja $z=-i$.
    
    \item Mitkä seuraavista väiteistä pätee ja mitkä ei? Polynomi $z^4+4$ voidaan esitt\"a\"a kahden polynomin tulona seuraavasti: 
    \begin{enumerate}
        \item $(z^2-2z+2)(z^2+2z+2)$ 
        \item $(z^2+2)^2$ 
        \item $(z^2+2)(z^2-2)$
    \end{enumerate}

    \textbf{Vastaus:} Ensimmäinen väite pätee. Kun toiseen väitteeseen soveltaa binomin neliön kaavaa, ilmestyy ylimääräinen termi $4z^2$, jota alkuperäisessä polynomissa ei ole. Kun kolmanteen väitteeseen käyttää summan ja erotuksen tulon muistikaavaa, tulee laskuun väärä etumerkki. 
    
\end{enumerate}
\end{enumerate}
\end{document}

\end{document}